\medskip\noindent\textbf{命题 2.4.1。}\phantomsection \label{sga3:VIIB:2.4.1}{\itshape 设 \(f:G\to K\) 为形式 \(k\)-群态射。若 \[
  H=\ker(f)
\] 在 \(k\) 上拓扑平坦，则 \(f\) 所诱导的同态 \[
  f':G/H\longrightarrow K
\] 为单态射。}

这是第四讲的结果的推论，\footnote{\label{sga3:VIIB:nde:93}见 \SGARefLink{sga3:iv:remark:5:2:6}{第四讲， 5.2.6}。}不过我们给出直接证明。设 \(T\) 为 \(k\) 上的有限长度形式簇，\(t\in(G/H)(T)\) 满足 \(f'\circ t\) 为 \(K(T)\) 的单位元。须证明 \(t\) 为 \((G/H)(T)\) 的单位元。设 \[
  p:G\longrightarrow G/H
\] 为投影，并设 \[
  X=T\times_{G/H}G.
\]。

由 \SGARefLink{sga3:VIIB:2.4}{2.4}，\(p\) 满射且拓扑平坦。因此 \[
  \operatorname{pr}_1:X\longrightarrow T.
\] 也如此。由 \SGARefLink{sga3:VIIB:1.3.1}{命题 1.3.1}，\(\operatorname{pr}_1\) 因而是满态射，故只需证明 \(t\circ\operatorname{pr}_1\) 为 \((G/H)(X)\) 的单位元。

设 \(\operatorname{pr}_2:X\to G\) 为第二投影。有 \[
  t\circ\operatorname{pr}_1
  =
  p\circ\operatorname{pr}_2,
\]，从而 \[
  1
  =
  f'\circ t\circ\operatorname{pr}_1
  =
  f'\circ p\circ\operatorname{pr}_2
  =
  f\circ\operatorname{pr}_2.
\]。因此正合列\begin{center}\phantomsection \label{sga3:VIIB:2.4.1:exact-sequence} % SGA3-VII-DIAG-112
% Authority: Exp7B-13oct24.pdf, local p. 54.
% Crop: figures/Exp7B_localp054_proposition241_exact_sequence.png
% Locator: tex/components/45_expose_VIIB_proposition241_kernel_quotient_image_en.tex, Loop 1 line 70.
\(
1\longrightarrow H\longrightarrow G\xrightarrow{\ f\ }K
\)
\end{center}表明 \(\operatorname{pr}_2\) 通过 \(H\) 分解。故 \(p\circ\operatorname{pr}_2\) 为零态射。由于 \[
  p\circ\operatorname{pr}_2
  =
  t\circ\operatorname{pr}_1,
\]，命题得证。

由此得到以下推论。以 \(\mathcal O(G)\)、\(\mathcal O(K)\) 和 \(\mathcal O(G/H)\) 分别表示 \(G\)、\(K\) 和 \(G/H\) 的仿射代数。已在 \SGARefLink{sga3:VIIB:2.4}{2.4} 中看到，\(p\) 诱导单射 \[
  \mathcal O(G/H)\lhook\joinrel\longrightarrow\mathcal O(G).
\]。此外，由于 \(f':G/H\to K\) 为单态射，\SGARefLink{sga3:VIIB:1.3.prop}{命题 1.3} 表明 \[
  \mathcal O(K)\longrightarrow\mathcal O(G/H)
\] 满射。因此：

\medskip\noindent\textbf{推论。}\phantomsection \label{sga3:VIIB:2.4.1:corollary}{\itshape 设 \(f:G\to K\) 为形式 \(k\)-群态射，且 \(H=\ker(f)\)。若 \(H\) 在 \(k\) 上拓扑平坦，则 \[
  \mathcal O(G/H)
\] 是 \(\mathcal O(K)\) 在 \(\mathcal O(G)\) 中的像。}
